\documentclass[11pt]{article}
\usepackage[margin=2.5cm]{geometry}
\usepackage[T1]{fontenc}
\usepackage[utf8]{inputenc}
\usepackage[slovene]{babel}
\usepackage{amsmath,amssymb,amsfonts}
\usepackage{array}
\usepackage{booktabs}
\usepackage{tabularx}

\begin{document}

\section*{Preizkus znanja: Racionalne funkcije}

\textbf{Ime in priimek:} \rule{6cm}{0.4pt}  \\
\textbf{Datum:} \rule{3cm}{0.4pt}  \hfill \textbf{Čas pisanja:} 60 minut \\
\textbf{Skupaj:} 40 točk  \hfill \textbf{Pripomočki:} pisalo in ravnilo; kalkulator ni potreben

\vspace{0.5cm}

\section*{Navodila}

\begin{itemize}
  \item Reši vse naloge in pri vsaki pokaži postopek.
  \item Pri racionalnih izrazih najprej zapiši pogoje oziroma izločene vrednosti.
  \item Rešitve enačb preveri v začetnem izrazu. Pri neenačbah izločene vrednosti ne sodijo v rešitev.
  \item Rezultate zapiši natančno, kjer je to mogoče; pri besedilni nalogi dodaj enoto.
\end{itemize}

\begin{tabular}{|l|c|c|c|c|c|c|c|}
\hline
Naloga & 1 & 2 & 3 & 4 & 5 & 6 & Skupaj \\
\hline
Točke & 6 & 5 & 6 & 8 & 7 & 8 & 40 \\
\hline
Čas (min) & 8 & 7 & 10 & 13 & 10 & 12 & 60 \\
\hline
\end{tabular}

\section*{Naloge}

\subsection*{1. Definicijsko območje in krajšanje -- 6 točk, 8 minut}

Dana je funkcija
\[
 f(x)=\frac{x^2-4}{x^2-x-2}.
\]

\textbf{a)} Razstavi števec in imenovalec na faktorje. (2 t) \\
\textbf{b)} Določi definicijsko območje funkcije. (2 t) \\
\textbf{c)} Funkcijo poenostavi in določi koordinati morebitne luknje na grafu. (2 t)

\subsection*{2. Racionalna enačba -- 5 točk, 7 minut}

Reši enačbo in zapiši izločeno vrednost:
\[
 \frac{2x+1}{x-1}=3.
\]

\subsection*{3. Racionalna neenačba -- 6 točk, 10 minut}

Reši neenačbo. Zapiši kritični vrednosti, predznake po intervalih in končno rešitev v intervalskem zapisu:
\[
 \frac{x-2}{x+3}<0.
\]

\subsection*{4. Lastnosti racionalne funkcije -- 8 točk, 13 minut}

Dana je funkcija
\[
 h(x)=\frac{3x^2+x-2}{x-1}.
\]

\textbf{a)} Določi definicijsko območje. (1 t) \\
\textbf{b)} Določi ničle in presečišče grafa z osjo $y$. (3 t) \\
\textbf{c)} Določi navpično in poševno asimptoto. (3 t) \\
\textbf{d)} Nariši okvirno skico ter označi dobljene ničle, presečišče in asimptoti. (1 t)

\subsection*{5. Racionalna enačba z več ulomki -- 7 točk, 10 minut}

Reši enačbo. Zapiši pogoje in preveri dobljeni rešitvi:
\[
 \frac{1}{x-2}+\frac{1}{x+2}=1.
\]

\subsection*{6. Uporaba racionalne enačbe -- 8 točk, 12 minut}

Stroj A sam izdela naročilo v 6 urah. Stroj A in stroj B, ki delujeta hkrati, naročilo izdelata v 4 urah. Predpostavi, da oba stroja delata z enakomerno hitrostjo.

Koliko časa bi za izdelavo istega naročila potreboval stroj B sam? Zapiši enačbo, pokaži postopek in odgovor podaj z enoto.

\end{document}
